% -------------------------------------------------------
%  Abstract
% -------------------------------------------------------

\begin{center}
	\textbf{چکیده}
\end{center}
\noindent
در این پروژه به بررسی و حل مشکل افت عملکرد ناشی از نقصان‌های برخورد در حافظه‌های نهان با روش نگاشت مستقیم پرداخته شده است. تداخل آدرس‌ها در این ساختار منجر به پدیدهٔ تلاطم و مراجعه‌های مکرر و پرهزینه به حافظهٔ اصلی می‌شود. برای رفع این چالش، یک حافظهٔ نهان قربانی تمام‌انجمنی طراحی و بین سطح اول حافظهٔ نهان و سطح پایین‌تر قرار داده شد. این معماری در شبیه‌ساز آکیتا پیاده‌سازی شده و عملکرد آن با استفاده از یک برنامهٔ محک با دسترسی‌های متناوب مورد ارزیابی قرار گرفت. نتایج شبیه‌سازی و بررسی شمارنده‌های سخت‌افزاری نشان می‌دهد که افزودن این پیمانه با مدیریت صحیح جابه‌جایی دوطرفه، ترافیک حافظهٔ اصلی را به شکل چشمگیری کاهش داده و زمان دسترسی میانگین به حافظه را بهبود می‌بخشد.

\vspace{\baselineskip}

\noindent\textbf{کلیدواژه‌ها:}
حافظهٔ نهان قربانی، نقصان برخورد، نگاشت مستقیم، تلاطم، زمان دسترسی میانگین.

\newpage